% problem-000078 2 分光光度分析
\section*{2. 分光光度分析}
在Mn(II)的催化下，某有机染料(CR)溶液被氧化剂(OX)氧化褪⾊。在⼀定的pH条件下该反应具有较⾼的灵敏度和选择性，据此可建⽴催化动⼒学光度法以测定⽔样中的痕量锰。催化反应⽰意如下：

$
\mathrm{CR} (\text {蓝色}) + \mathrm{OX} \xrightarrow {\mathrm{Mn(II)}} \text {褪色产物}
$

CR褪⾊反应的速率⽅程为：

$
r = - \frac {\mathrm{d} c _ {\mathrm{CR}}}{\mathrm{d} t} = k c _ {C R} ^ {\alpha} c _ {O X} ^ {\beta} c _ {\mathrm{Mn(II)}} ^ {\gamma}
$

已知体系中⽆Mn(II)时褪⾊反应不发⽣。反应过程中，氧化剂⼤⼤过量。

\subsection*{2-1}
固定溶液中Mn(II)的浓度，改变CR溶液的初始浓度，测得溶液在6.0min和9.0min时的吸光度，部分数据⻅表2-1，推求α的值。

表 2-1

\textit{（原卷此处含表格/仪器试剂表，详见 PDF 或 MD 源文件。）}

\subsection*{2-2}
实验表明， $\gamma = \mathrm { l _ { c } }$ 。推出固定时间间隔Δ�时，相邻两时刻的吸光度⽐值与Mn(II)浓度之间的定量关系式。

\subsection*{2-3}
固定 CR 起始浓度及 Mn(II)的浓度，分别测定了 363.15 K 和 373.15 K 下体系在 $t = 8 . 0 \operatorname* { m i n }$ 和 10.0min时的吸光度。实验结果⻅表2-2，据此求算反应的表观活化能。

表 2-2

\textit{（原卷此处含表格/仪器试剂表，详见 PDF 或 MD 源文件。）}
